% Generated draft: requires mathematical and source review.
\documentclass[11pt]{article}
\usepackage[margin=25mm]{geometry}
\usepackage{fontspec}
\setmainfont{Latin Modern Roman}
\usepackage{amsmath,amssymb,mathtools,unicode-math}
\setmathfont{Latin Modern Math}
\usepackage{enumitem,needspace,xurl,hyperref}
\setlength{\parindent}{0pt}
\setlength{\parskip}{5pt}
\setlength{\emergencystretch}{3em}
\newcommand{\N}{\mathbb{N}}
\newcommand{\Z}{\mathbb{Z}}
\newcommand{\Q}{\mathbb{Q}}
\newcommand{\R}{\mathbb{R}}
\newcommand{\C}{\mathbb{C}}
\newcommand{\F}{\mathbb{F}}
\newcommand{\A}{\mathbb{A}}
\newcommand{\PP}{\mathbb{P}}
\newcommand{\E}{\mathbb{E}}
\newcommand{\eps}{\varepsilon}
\newcommand{\dd}{\,\mathrm{d}}
\begin{document}
\section*{UnsolvedMath ID: 800685}
\textbf{Source problem code:} catalogue entry
\subsection{Definitions and mathematical statement}
The supplied record uses a space of maps $\mathscr{H}$, a translation operator on $\mathscr{H}$, and a subset $\mathscr{U}\subseteq\mathscr{H}$ without defining them. Consequently, the following formulation is conditional on specifying these objects and their measurable structure.

Write $T:\mathscr{H}\to\mathscr{H}$ for the translation operator referred to in the source, and let $\Sigma$ denote the intended sigma-algebra on $\mathscr{H}$. Assuming $T$ is measurable, a probability measure $\mu$ on $(\mathscr{H},\Sigma)$ is invariant if
\[
\mu(T^{-1}A)=\mu(A)\qquad\text{for every }A\in\Sigma.
\]
Let $\mathscr{C}\subseteq\mathscr{H}$ denote the set of constant maps. Assuming $\mathscr{C}\in\Sigma$, the source calls $\mu$ \emph{non-trivial} when $\mu(\mathscr{C})=0$. An invariant probability measure is ergodic if every invariant measurable set has measure zero or one; the precise convention for invariance must agree with the intended translation action.

\textbf{Problem (Dinh--Sibony, 2018, Problem 12).} Given a non-trivial invariant probability measure $\mu$ on $\mathscr{H}$ as in the source, what can be said about
\[
\mu(\mathscr{U})?
\]
This expression requires $\mathscr{U}$ to be measurable for the intended measure. The record does not define $\mathscr{U}$ or specify a proposed value or bound. It also does not make clear whether the question concerns all such invariant measures or specifically the ergodic measures discussed immediately before it. These missing qualifications prevent a fully specified mathematical statement.
\subsection{Short English statement}
Dinh and Sibony ask what can be said about the measure of $\mathscr{U}$ under a non-trivial translation-invariant probability measure on $\mathscr{H}$. The supplied excerpt omits the definitions of the space, the translation operator, and the subset, and leaves the intended ergodicity assumption unclear.
\subsection{Sources}
\begin{itemize}
\item Supplied local UnsolvedMath registry record: permanent ID $800685$, source problem code \texttt{catalogue entry}, title ``Holomorphic Foliation Theory --- Problem 12.'' Attribution: Dinh, Sibony; proposed year: 2018. The registry labels the problem open; this label is supplied metadata, not an independently verified status assessment.
\item Dinh, Sibony, \emph{Some Open Problems on Holomorphic Foliation Theory} (2018), Problem 12. Supplied source URL: \url{https://arxiv.org/abs/1807.05640}. The source record reports access on 2026-10-16 and extraction by \texttt{arxiv\_latex}; these provenance details are supplied and have not been independently verified.
\item The excerpt attributes existence of non-trivial invariant probability measures to Benjamin Weiss, under the reference label [Weiss], and says these measures are studied by Buhovsky, Gl\"ucksam, Logunov and Sodin, under [BGLS]. Full bibliographic details for these references are not supplied. The excerpt further attributes the existence of non-trivial ergodic measures to standard results from ergodic theory, without specifying a theorem or citation.
\end{itemize}
\end{document}
